\section{核心代码}
% 写作提示：附录应放与正文一致、完整可运行且能复现关键结果的真实代码，并说明环境、依赖、输入和输出。
% 下方代码仅为 listings 排版示例；正式论文应替换为实际程序，较长代码可用 \lstinputlisting 引入文件。
% 如实际使用 AI 辅助编程，须在程序文件开头保留并如实填写来源注释；未使用则删除示例注释。
\begin{Python}{核心求解程序}
# 本程序及代码是在人工智能工具辅助下完成的。
# 工具名称：[实际使用的工具]；版本/型号：[版本]
# 开发机构/公司：[机构]；版本发布日期：[YYYY-MM-DD]
import numpy as np
import pandas as pd
from sklearn.ensemble import RandomForestRegressor
from sklearn.model_selection import train_test_split
from sklearn.metrics import r2_score

data = pd.read_csv('data.csv')
X = data.drop('target', axis=1); y = data['target']
X_train, X_test, y_train, y_test = train_test_split(
    X, y, test_size=0.2, random_state=42
)
model = RandomForestRegressor(n_estimators=100, random_state=42)
model.fit(X_train, y_train)
print(f'R2: {r2_score(y_test, model.predict(X_test)):.4f}')
\end{Python}
